% ============================================================================
%  A Field Theory of Sentiment
%  Affect as a circulating, interoceptive resource-allocation field that
%  modulates variance-minimizing computation.
%
%  Self-contained: developed from first principles; cites only the standard
%  external literature.
% ============================================================================
\documentclass[11pt,a4paper]{article}

\usepackage[utf8]{inputenc}
\usepackage[T1]{fontenc}
\usepackage{lmodern}
\usepackage[a4paper,margin=2.4cm]{geometry}
\usepackage{amsmath,amssymb,amsthm,mathtools}
\usepackage{enumitem}
\usepackage{booktabs}
\usepackage{array}
\usepackage{xcolor}
\usepackage{microtype}
\usepackage[numbers,sort&compress]{natbib}
\usepackage[colorlinks=true,linkcolor=blue!45!black,citecolor=blue!45!black,
            urlcolor=blue!55!black]{hyperref}
\usepackage{cleveref}

% ---- theorem environments ----
\theoremstyle{plain}
\newtheorem{theorem}{Theorem}[section]
\newtheorem{lemma}[theorem]{Lemma}
\newtheorem{proposition}[theorem]{Proposition}
\newtheorem{corollary}[theorem]{Corollary}
\theoremstyle{definition}
\newtheorem{definition}[theorem]{Definition}
\newtheorem{assumption}[theorem]{Assumption}
\newtheorem{example}[theorem]{Example}
\theoremstyle{remark}
\newtheorem{remark}[theorem]{Remark}
\newtheorem{principle}[theorem]{Principle}

% ---- operators and notation ----
\DeclareMathOperator{\dom}{dom}
\DeclareMathOperator{\supp}{supp}
\DeclareMathOperator{\Tr}{Tr}
\DeclareMathOperator{\Dg}{Dg}
\DeclareMathOperator*{\argmin}{arg\,min}
\DeclareMathOperator*{\argmax}{arg\,max}
\newcommand{\R}{\mathbb{R}}
\newcommand{\Rp}{\mathbb{R}_{\ge 0}}
\newcommand{\N}{\mathbb{N}}
\newcommand{\Unit}{[0,1]}
\newcommand{\simplex}{\Delta}                 % probability simplex
\newcommand{\field}{\mathbf{S}}               % sentiment field (vector)
\newcommand{\alloc}{\mathbf{a}}               % allocation
\newcommand{\feat}{\mathbf{g}}                % feature/affordance map
\newcommand{\Pot}{\Phi}                       % modulated potential
\newcommand{\Free}{\mathcal{F}}               % free energy
\newcommand{\coh}{R}                           % coherence order parameter
\newcommand{\KL}{D_{\mathrm{KL}}}
\newcommand{\inner}[2]{\langle #1,\,#2\rangle}
\newcommand{\norm}[1]{\lVert #1\rVert}
\newcommand{\dd}{\,\mathrm{d}}
\newcommand{\setpt}{\field^{\!\star}}

\title{\textbf{A Field Theory of Sentiment}\\[3pt]
\large Affect as a Circulating, Interoceptive Resource-Allocation Field\\
that Modulates Variance-Minimizing Computation}

\author{Kundai Farai Sachikonye\\[2pt]
\small Technical University of Munich\\
\small \texttt{kundai.sachikonye@wzw.tum.de}}

\date{\today}

\begin{document}
\maketitle

\begin{abstract}
\noindent
An autonomous computational system that pursues an open-ended goal must, at
every instant, allocate finite resources among many candidate sub-computations,
none of which individually apprehends the whole goal. We propose that the
control variable governing this allocation is best modelled not as a scalar
priority but as a \emph{field}: a circulating, self-sensing quantity that
modulates which candidate computations become stable and therefore proceed. We
call this quantity \emph{sentiment}, by analogy with the role that affect plays
in biological cognition, where it is increasingly understood as an
\emph{interoceptive} read-out of the organism's internal economy rather than an
overlay on neutral thought. We give a complete, self-contained mathematical
theory of such a field. The theory has four interlocking components. (i) A
\emph{modulation} principle: candidate computations are local minima of a
potential $\Pot_\field(x)=V(x)-\inner{\field}{\feat(x)}$ in which the field
$\field$ tilts an intrinsic cost (a ``variance'') by the affordances $\feat$ of
each candidate; the set of stable candidates therefore moves smoothly with the
field except at bifurcations, which we identify with affective regime change.
(ii) A \emph{circulation} law: a conserved resource is transported across the
candidate population and periodically re-injected by a pump, so that allocation
is the flow of a conserved budget rather than an unconstrained gain. (iii) An
\emph{interoception} law: the field is driven by the system's own observed
coherence and residual error through a homeostatic feedback, so that the field
reports the internal state it also governs. (iv) An \emph{economy}: given the
field, the budget is allocated by a maximum-entropy (Gibbs) rule that is the
unique minimizer of a free energy, equivalently the stationary distribution of
an entropy-regularised replicator dynamics, with the field acting as the price
(Lagrange multiplier) of the budget constraint. We prove the theory's basic
guarantees: budget conservation, boundedness, smooth modulation, convergence of
the allocation to the Gibbs minimizer, a coherence phase transition at a
critical coupling, and stability of the coupled slow-field/fast-allocation
system under a timescale-separation condition. We show that \emph{memory} is the
path integral of the field, $M(t)=\field(t)-\field(0)$, making field and memory
a derivative--integral pair. Finally we isolate the theory's interface to a
host runtime: the sentiment field induces a map from system state to an
allocation over candidate computations---a scheduler priority that is
self-tuning, conservative, and accountable. Throughout, the construction is
related to the standard literature on dynamical systems, feedback control,
evolutionary dynamics, convex optimization, synchronization, and the
interoceptive theory of affect.

\medskip
\noindent\textbf{Keywords:} sentiment field, resource allocation, replicator
dynamics, maximum entropy, homeostasis, interoception, synchronization,
variational principle, gradient flow, attention economy.
\end{abstract}

\tableofcontents
\bigskip\hrule

% ===========================================================================
\section{Introduction}
\label{sec:intro}
% ===========================================================================

\subsection{The allocation problem of open-ended computation}

Consider a system charged with an open-ended task: to investigate a phenomenon,
to design an artefact, to reach a goal whose full structure is not known in
advance. Such a system does not execute a fixed plan. Instead it maintains, at
each moment, a \emph{population} of candidate sub-computations---partial
hypotheses, tentative subgoals, lines of inquiry---and it must decide which of
them to pursue with its finite resources of time, memory, and computation. The
defining difficulty is that no single candidate apprehends the whole task: each
is local, partial, and uncertain. The system's competence lies not in any one
candidate but in how it \emph{allocates} effort across them and reallocates as
the situation changes.

The naive control variable for such allocation is a scalar priority assigned to
each candidate. We argue that this is inadequate in three respects, each of
which a richer construct must supply. First, priority must be \emph{coupled}:
the value of pursuing one candidate depends on the state of the others and on
the global situation, so the control variable should be a field over the task
space, not a list of independent numbers. Second, priority must be
\emph{conserved}: resources are finite, so raising the priority of one candidate
must lower that of another; allocation is the redistribution of a conserved
budget, not an unconstrained gain. Third, priority must be \emph{self-sensing}:
the appropriate allocation depends on how the computation is itself faring---on
whether it is coherent or thrashing, near its goal or far---so the control
variable must read the system's own internal state, which it also governs.

\subsection{Sentiment as an interoceptive field}

We name the control variable that meets these three requirements
\emph{sentiment}. The name is chosen deliberately. In biological cognition,
affect was long modelled as an overlay---a colouring of otherwise neutral
representations. The modern view is the reverse: affect is an
\emph{interoceptive} signal, a read-out of the body's internal economy
(metabolic, autonomic, immune), that the brain uses to allocate its own
resources and to bias which thoughts become salient and stable
\citep{craig2002,damasio1994,barrett2017}. Emotion is, on this view, not
decoration but the very field that determines which lines of processing
stabilise under the current internal conditions. Our theory abstracts exactly
this role. We do not claim to model biological affect; we claim that the
\emph{computational} role affect is now understood to play---a circulating,
self-sensing field that modulates resource allocation across competing
sub-computations---is a precise and useful control architecture, and we give it
a precise mathematics.

The construction rests on a circulation metaphor made literal. A conserved
resource---call it the \emph{budget}---is pumped through the population of
candidates, delivering both sustenance (the computational effort a candidate
needs to remain viable) and information (the field value that biases it). This
is the same dual role that a circulatory medium plays in a living system, where
the blood simultaneously sustains tissue and reports its condition. The pump
imposes a rhythm; the flow obeys a conservation law; the field that the flow
carries is set by the system's interoceptive sensing of its own coherence.

\subsection{The four laws}

The theory is organised around four laws, each a section of this paper.

\begin{enumerate}[leftmargin=1.5em]
\item \textbf{Modulation} (\cref{sec:modulation}). Candidate computations are
the stable equilibria---local minima---of a potential
$\Pot_\field(x)=V(x)-\inner{\field}{\feat(x)}$ in which an intrinsic cost $V$
(the ``variance'' a candidate seeks to minimise) is tilted by the field
$\field$ along the candidate's affordances $\feat(x)$. The stable set moves
smoothly with $\field$ except at bifurcations, which are the formal signature of
a change of affective regime.

\item \textbf{Circulation} (\cref{sec:circulation}). The budget is a conserved
density transported across candidates and re-injected by a pump; allocation is
its flow, and total budget is invariant.

\item \textbf{Interoception} (\cref{sec:interoception}). The field is driven by
the system's own coherence and residual error through homeostatic feedback; it
reports the state it governs, closing the loop.

\item \textbf{Economy} (\cref{sec:economy}). Given the field, the budget is
allocated by a maximum-entropy rule that uniquely minimises a free energy and is
the stationary point of an entropy-regularised replicator dynamics; the field is
the price of the budget constraint.
\end{enumerate}

\Cref{sec:memory} identifies memory with the path integral of the field;
\cref{sec:metatheory} collects the metatheory (conservation, boundedness,
convergence, the phase transition, slow--fast stability); \cref{sec:interface}
states the interface to a host runtime; \cref{sec:example} gives a worked
example; \cref{sec:related} situates the theory; \cref{sec:conclusion}
concludes.

\subsection{Notation}
$\R,\Rp,\N$ are the reals, nonnegative reals, naturals; $\Unit$ the unit
interval. $\simplex^{N-1}=\{\alloc\in\Rp^N:\sum_i a_i=1\}$ is the probability
simplex and $\mathrm{int}\,\simplex^{N-1}$ its interior. $\inner{\cdot}{\cdot}$
is the Euclidean inner product, $\norm{\cdot}$ its norm. For a smooth
$h\colon\R^n\to\R$, $\nabla h$ and $\nabla^2 h$ are gradient and Hessian; for
$\feat\colon\R^n\to\R^d$, $D\feat$ is the Jacobian. $H(\alloc)=-\sum_i a_i\ln
a_i$ is the Shannon entropy \citep{shannon1948} and $\KL(\alloc\,\|\,\alloc')=
\sum_i a_i\ln(a_i/a'_i)$ the Kullback--Leibler divergence \citep{cover2006}.

% ===========================================================================
\section{The Modulation Law}
\label{sec:modulation}
% ===========================================================================

\subsection{Candidate computations and the intrinsic cost}

\begin{definition}[Candidate population]
\label{def:population}
A \emph{candidate population} at an instant is a finite set
$C=\{x_1,\dots,x_N\}\subset\mathcal{X}$ of points in a base space
$\mathcal{X}\subseteq\R^n$, the \emph{task space}. Each $x_i$ is a candidate
sub-computation, identified with its location in feature coordinates. Two maps
are given on $\mathcal{X}$: an \emph{intrinsic cost} $V\colon\mathcal{X}\to\Rp$,
smooth and bounded below, which each candidate locally seeks to decrease (its
``variance''); and an \emph{affordance map} $\feat\colon\mathcal{X}\to\R^d$,
smooth, whose components $g_k(x)$ measure the extent to which $x$ exhibits the
$k$-th salient feature.
\end{definition}

The intrinsic cost is the language-agnostic stand-in for whatever local
objective a candidate minimises---residual uncertainty, prediction error,
distance to a subgoal. The theory does not depend on its identity, only on its
smoothness and lower-boundedness.

\subsection{The sentiment field and the modulated potential}

\begin{definition}[Sentiment field]
\label{def:field}
A \emph{sentiment field} is a (time-dependent) vector $\field\in\R^d$, the value
of the affective field at the population's current locus. (Globally $\field$ is
a field $\mathcal{X}\to\R^d$; throughout we work with its value on the active
region and write $\field$ for $\field(\bar x)$ at the population centroid
$\bar x$.) The field acts on candidates through the \emph{modulated potential}
\begin{equation}
\Pot_\field(x)\;=\;V(x)\;-\;\inner{\field}{\feat(x)}.
\label{eq:potential}
\end{equation}
\end{definition}

\begin{principle}[Modulation]
\label{prin:modulation}
A candidate $x$ is \emph{stable under the field $\field$} iff it is a local
minimum of $\Pot_\field$. The field does not select candidates by fiat; it
reshapes the landscape on which candidates settle. Raising $S_k$ rewards
affordance $g_k$, lowering the modulated cost of candidates rich in feature $k$
and so making them stable; lowering $S_k$ suppresses them.
\end{principle}

This is the formal content of the slogan that sentiment determines \emph{which
thoughts can stabilise}: it is the field that, by tilting the variance
landscape, decides which local minima exist and where they sit.

\subsection{Smooth modulation and regime change}

\begin{theorem}[Smooth modulation]
\label{thm:modulation}
Let $x^\ast$ be a nondegenerate local minimum of $\Pot_{\field_0}$, i.e.\
$\nabla\Pot_{\field_0}(x^\ast)=0$ and the Hessian
\[
H(x^\ast,\field_0)\;=\;\nabla^2 V(x^\ast)-\sum_{k=1}^{d}(\field_0)_k
\nabla^2 g_k(x^\ast)\;\succ\;0 .
\]
Then there is a neighbourhood $U\ni\field_0$ and a unique smooth branch
$x^\ast\colon U\to\mathcal{X}$ with $x^\ast(\field_0)=x^\ast$, each
$x^\ast(\field)$ a nondegenerate minimum of $\Pot_\field$, and
\begin{equation}
\frac{\partial x^\ast}{\partial\field}\;=\;H(x^\ast,\field)^{-1}\,
D\feat\big(x^\ast\big)^{\!\top}.
\label{eq:sensitivity}
\end{equation}
\end{theorem}
\begin{proof}
Apply the implicit function theorem to $F(x,\field)=\nabla\Pot_\field(x)=
\nabla V(x)-D\feat(x)^{\top}\field$. We have $F(x^\ast,\field_0)=0$ and
$\partial_x F(x^\ast,\field_0)=H(x^\ast,\field_0)$, which is invertible by
nondegeneracy. Hence a unique smooth $x^\ast(\field)$ solves $F=0$ near
$\field_0$, and positive-definiteness persists on a neighbourhood by continuity
of eigenvalues, so each $x^\ast(\field)$ is a minimum. Differentiating
$F(x^\ast(\field),\field)=0$ gives $H\,\partial_\field x^\ast - D\feat^{\top}=0$,
which is \cref{eq:sensitivity}.
\end{proof}

\begin{definition}[Affective regime change]
\label{def:bifurcation}
An \emph{affective regime change} is a value $\field_c$ at which
$H(x^\ast,\field_c)$ becomes singular: a stable candidate is created or
destroyed, or two merge. By \cref{thm:modulation} this is exactly where the
smooth branch fails; it is a bifurcation of $\Pot_\field$
\citep{strogatz2015,hirsch2013}.
\end{definition}

\begin{remark}[Why the stable set is the right object]
Modulation makes precise that the field's effect is \emph{structural}: between
regime changes the stable candidates track the field smoothly
(\cref{eq:sensitivity}); at regime changes the very inventory of viable
candidates changes. A controller that manipulates $\field$ therefore exercises
two qualitatively different kinds of influence---continuous steering within a
regime, and discontinuous reconfiguration across regimes---without ever
selecting a candidate directly.
\end{remark}

% ===========================================================================
\section{The Circulation Law}
\label{sec:circulation}
% ===========================================================================

Allocation is the distribution of a conserved budget across the population. We
make the conservation explicit and give the budget a transport dynamics with a
pump, so that the field's modulation acts on a flow rather than on free gains.

\begin{definition}[Budget and allocation]
\label{def:budget}
The \emph{budget} is a fixed total $B>0$ of an abstract resource (effort,
compute, attention). An \emph{allocation} is $\alloc\in\simplex^{N-1}$, where
$a_i$ is the fraction of $B$ delivered to candidate $x_i$; the resource
delivered is $B\alloc$. We model $\alloc$ as a density that evolves in time.
\end{definition}

\begin{definition}[Transport with a pump]
\label{def:transport}
Let $\mu_i$ be the \emph{marginal value} of resource at candidate $i$
(\cref{sec:economy} fixes $\mu_i=f_i$). The allocation obeys a discrete
continuity equation with a conservative flux and a pump source:
\begin{equation}
\dot a_i \;=\; \underbrace{\sum_{j\sim i}J_{j\to i}}_{\text{conservative flux}}
\;+\; \underbrace{\sigma_i(t)}_{\text{pump}},
\qquad J_{j\to i}=-J_{i\to j},\qquad \sum_i\sigma_i(t)=0,
\label{eq:continuity}
\end{equation}
where $j\sim i$ ranges over a coupling graph on the population, the flux is
antisymmetric (resource leaving $j$ arrives at $i$), and the pump $\sigma$
redistributes without creating or destroying budget (e.g.\ a periodic
re-seeding that injects exploration while drawing equally from the incumbents).
\end{definition}

\begin{theorem}[Budget conservation]
\label{thm:conservation}
Under \cref{eq:continuity}, $\sum_i a_i$ is invariant; if $\alloc(0)\in
\simplex^{N-1}$ then $\alloc(t)\in\simplex^{N-1}$ for all $t$, and the delivered
resource $B\sum_i a_i=B$ is conserved.
\end{theorem}
\begin{proof}
$\frac{\dd}{\dd t}\sum_i a_i=\sum_i\sum_{j\sim i}J_{j\to i}+\sum_i\sigma_i$. The
first term is $\sum_{\{i,j\}}(J_{j\to i}+J_{i\to j})=0$ by antisymmetry; the
second is $0$ by the pump constraint. Hence $\sum_i a_i$ is constant; with
initial sum $1$ it stays $1$. Nonnegativity is preserved provided the flux and
pump vanish at the boundary $a_i=0$ (no outflow from an empty candidate), which
we assume of $J,\sigma$.
\end{proof}

\begin{remark}[The pump as rhythm and exploration]
The pump $\sigma$ is the theory's ``heart''. A natural choice is a slow
oscillation $\sigma_i(t)=\epsilon\,\omega(t)\,(\pi_i-a_i)$ with
$\sum_i\pi_i=1$ a re-seeding distribution and $\omega$ periodic: it gently and
periodically pulls the allocation toward $\pi$ (exploration) without violating
conservation, since $\sum_i\sigma_i=\epsilon\omega(\sum_i\pi_i-\sum_i a_i)=0$ on
the simplex. The oscillation prevents the allocation from collapsing onto a
single incumbent and is the mechanism by which a system keeps ``circulating''
rather than fixating.
\end{remark}

% ===========================================================================
\section{The Interoception Law}
\label{sec:interoception}
% ===========================================================================

The field is not exogenous: it is set by the system's reading of its own state.
We give two interoceptive signals---a \emph{coherence} and a \emph{residual
error}---and a homeostatic law driving the field toward a setpoint determined by
them \citep{cannon1929,ashby1960,astrom2008}.

\begin{definition}[Coherence and residual error]
\label{def:coherence}
Assign to each candidate a \emph{phase} $u(x_i)\in\mathbb{S}^{m-1}$ (a unit
vector encoding the direction the candidate ``points'', e.g.\ its dominant
feature). The \emph{coherence} of an allocation is the order parameter
\begin{equation}
\coh(\alloc)\;=\;\Big\lVert\,\sum_i a_i\,u(x_i)\,\Big\rVert\;\in\;\Unit,
\label{eq:coherence}
\end{equation}
which is near $1$ when the allocated mass points one way (the population agrees)
and near $0$ when it is dispersed \citep{kuramoto1984,strogatz2000,winfree1980}.
The \emph{residual error} is the expected intrinsic cost
$\varepsilon(\alloc)=\sum_i a_i V(x_i)$.
\end{definition}

\begin{definition}[Interoceptive homeostasis]
\label{def:homeostasis}
The field relaxes toward a setpoint $\setpt(\coh,\varepsilon)$ that depends on
the interoceptive signals:
\begin{equation}
\dot\field\;=\;-\gamma\big(\field-\setpt(\coh(\alloc),\varepsilon(\alloc))\big),
\qquad \gamma>0.
\label{eq:homeostasis}
\end{equation}
The map $\setpt$ encodes the system's affective ``policy'': for instance,
$\setpt$ may raise the weight on exploratory affordances when coherence is high
and error is not improving (escape a confident dead end), and raise the weight
on consolidating affordances when coherence is low (gather the dispersed
population).
\end{definition}

\begin{proposition}[Homeostatic stability]
\label{prop:homeostasis}
For fixed interoceptive signals (hence fixed setpoint $\setpt$), the field
dynamics \cref{eq:homeostasis} is globally exponentially stable with rate
$\gamma$: $\norm{\field(t)-\setpt}\le e^{-\gamma t}\norm{\field(0)-\setpt}$.
\end{proposition}
\begin{proof}
$\frac{\dd}{\dd t}\tfrac12\norm{\field-\setpt}^2=\inner{\field-\setpt}
{\dot\field}=-\gamma\norm{\field-\setpt}^2$, so $\norm{\field-\setpt}^2$ decays
as $e^{-2\gamma t}$ by Grönwall; take square roots \citep{khalil2002}.
\end{proof}

\begin{remark}[Interoception in one line]
Equations \cref{eq:coherence}--\cref{eq:homeostasis} are the formal content of
``sentiment is interoception'': the field is a function of the system's own
coherence and error, i.e.\ of how the computation is going, and it feeds back to
reshape (via \cref{eq:potential}) which candidates are stable. The field thus
both \emph{reports} and \emph{regulates} the internal state, exactly as affect
is held to do \citep{craig2002,barrett2017}.
\end{remark}

% ===========================================================================
\section{The Economy}
\label{sec:economy}
% ===========================================================================

We now fix how, given the field, the budget is allocated. The principle is
\emph{maximum entropy under an expected-fitness constraint}: among all
allocations achieving a given expected modulated benefit, choose the least
committal \citep{jaynes1957}. This yields the Gibbs rule, which we show is the
unique minimiser of a free energy and the stationary point of an
entropy-regularised replicator dynamics, with the field appearing as the price
of the budget.

\begin{definition}[Fitness]
\label{def:fitness}
The \emph{fitness} of a candidate under field $\field$ is the negative modulated
cost evaluated at its stable location:
$f_i(\field)=-\Pot_\field(x_i)=-V(x_i)+\inner{\field}{\feat(x_i)}$. Higher
fitness means lower modulated cost, hence greater claim on the budget.
\end{definition}

\begin{definition}[Free energy and the Gibbs allocation]
\label{def:gibbs}
For inverse temperature $\beta>0$ define the \emph{free energy}
\begin{equation}
\Free_\beta(\alloc)\;=\;-\sum_i a_i f_i \;-\;\tfrac{1}{\beta}H(\alloc)
\;=\;-\inner{\alloc}{f}-\tfrac1\beta H(\alloc),\qquad \alloc\in\simplex^{N-1},
\label{eq:freeenergy}
\end{equation}
and the \emph{Gibbs allocation} $\alloc^\beta$ with
$a^\beta_i=\dfrac{e^{\beta f_i}}{\sum_j e^{\beta f_j}}$.
\end{definition}

\begin{theorem}[The Gibbs allocation is the unique free-energy minimiser]
\label{thm:gibbs}
$\Free_\beta$ is strictly convex on $\simplex^{N-1}$ and attains its unique
minimum at $\alloc^\beta$. Equivalently, $\alloc^\beta$ is the maximum-entropy
allocation subject to a fixed expected fitness $\inner{\alloc}{f}$, and the
field-induced terms $\inner{\field}{\feat(x_i)}$ enter $f_i$ linearly, so the
field acts as a vector of \emph{prices} on the affordances.
\end{theorem}
\begin{proof}
On $\simplex^{N-1}$, $-H$ is strictly convex (its Hessian is $\mathrm{diag}
(1/a_i)\succ0$) and $-\inner{\alloc}{f}$ is linear, so $\Free_\beta$ is strictly
convex and has a unique minimiser \citep{boyd2004}. Forming the Lagrangian
$L=\Free_\beta(\alloc)+\lambda(\sum_i a_i-1)$ and setting
$\partial L/\partial a_i=-f_i+\tfrac1\beta(\ln a_i+1)+\lambda=0$ gives
$a_i\propto e^{\beta f_i}$; normalisation yields $\alloc^\beta$. The
max-entropy characterisation is the standard duality: maximising $H(\alloc)$
subject to $\inner{\alloc}{f}=c$ and $\sum a_i=1$ has stationarity condition
$\ln a_i=\beta f_i+\text{const}$ with $\beta$ the multiplier on the fitness
constraint \citep{jaynes1957}. The affordance prices are the components of
$\field$ because $f_i=-V(x_i)+\sum_k S_k g_k(x_i)$.
\end{proof}

\begin{definition}[Entropy-regularised replicator dynamics]
\label{def:replicator}
The allocation evolves by the \emph{logit/replicator} flow
\begin{equation}
\dot a_i\;=\;a_i\big(f_i-\bar f\big)\;-\;\tfrac1\beta\,a_i\big(\ln a_i-
\overline{\ln a}\big),\qquad \bar f=\inner{\alloc}{f},\ \
\overline{\ln a}=\inner{\alloc}{\ln\alloc},
\label{eq:replicator}
\end{equation}
the replicator dynamics \citep{taylor1978,hofbauer1998,maynardsmith1982} with an
entropy-regularising drift. The unregularised flow ($\beta\to\infty$) is the
classical replicator equation; the drift keeps the allocation in the interior.
\end{definition}

\begin{theorem}[Convergence of the allocation]
\label{thm:repconv}
For fixed fitness $f$, the free energy $\Free_\beta$ is a strict Lyapunov
function for \cref{eq:replicator} on $\mathrm{int}\,\simplex^{N-1}$:
$\dot{\Free}_\beta\le0$ with equality only at $\alloc^\beta$. Hence every
interior trajectory converges to the Gibbs allocation $\alloc^\beta$, which is
globally asymptotically stable.
\end{theorem}
\begin{proof}
Write $f^{\mathrm{eff}}_i=f_i-\tfrac1\beta\ln a_i$, so \cref{eq:replicator} is
$\dot a_i=a_i(f^{\mathrm{eff}}_i-\overline{f^{\mathrm{eff}}})$ with
$\overline{f^{\mathrm{eff}}}=\inner{\alloc}{f^{\mathrm{eff}}}$, a replicator
flow in the effective fitness. Note $f^{\mathrm{eff}}_i=-\partial\Free_\beta/
\partial a_i+\text{const}$. A standard computation gives
$\dot{\Free}_\beta=\sum_i\frac{\partial\Free_\beta}{\partial a_i}\dot a_i
=-\big(\overline{(f^{\mathrm{eff}})^2}-(\overline{f^{\mathrm{eff}}})^2\big)
=-\mathrm{Var}_{\alloc}(f^{\mathrm{eff}})\le0$, the replicator analogue of
Fisher's fundamental theorem \citep{hofbauer1998}. Equality holds iff
$f^{\mathrm{eff}}$ is constant across the support, i.e.\ $\ln a_i=\beta f_i
+\text{const}$, i.e.\ $\alloc=\alloc^\beta$. Strict convexity of
$\Free_\beta$ (\cref{thm:gibbs}) makes $\alloc^\beta$ the unique minimum, so by
LaSalle's invariance principle \citep{khalil2002} every interior trajectory
converges to it.
\end{proof}

\begin{corollary}[Temperature controls commitment]
\label{cor:temperature}
As $\beta\to\infty$ the Gibbs allocation concentrates on
$\argmax_i f_i$ (full commitment to the fittest candidate); as $\beta\to0$ it
tends to the uniform allocation (maximal exploration). The inverse temperature
$\beta$ is the system's exploitation--exploration dial.
\end{corollary}
\begin{proof}
Immediate from $a^\beta_i\propto e^{\beta f_i}$: the softmax tends to the
indicator of the maximiser as $\beta\to\infty$ and to uniform as $\beta\to0$
\citep{cover2006}.
\end{proof}

% ===========================================================================
\section{Memory as the Field's Path Integral}
\label{sec:memory}
% ===========================================================================

The field carries its own history. We define memory as the accumulated change
of the field along the computation's trajectory and record its elementary
properties.

\begin{definition}[Affective memory]
\label{def:memory}
The \emph{affective memory} accumulated up to time $t$ is the path integral of
the field's rate of change,
\begin{equation}
M(t)\;=\;\int_0^t\dot\field(\tau)\,\dd\tau\;=\;\field(t)-\field(0).
\label{eq:memory}
\end{equation}
More finely, the \emph{traversed memory} $\widetilde M(t)=\int_0^t
\norm{\dot\field(\tau)}\,\dd\tau$ records the total variation of the field along
the path, distinguishing trajectories that reach the same field value by
different routes.
\end{definition}

\begin{proposition}[Field--memory duality and path dependence]
\label{prop:memory}
$M$ and $\field$ are a derivative--integral pair: $\dot M=\dot\field$ and
$M(t)=\field(t)-\field(0)$. The reduced memory $M$ depends only on the
endpoints, but the traversed memory $\widetilde M$ is path-dependent:
$\widetilde M(t)\ge\norm{M(t)}$, with equality iff the field moves
monotonically along a ray. Consequently two computations reaching the same field
state through different affective histories are distinguished by $\widetilde M$.
\end{proposition}
\begin{proof}
$\dot M=\dot\field$ and the endpoint formula are the fundamental theorem of
calculus. The inequality is the triangle inequality for the path-length
functional: $\widetilde M=\int_0^t\norm{\dot\field}\ge\norm{\int_0^t\dot\field}
=\norm{M}$, with equality iff $\dot\field$ keeps a fixed direction.
\end{proof}

\begin{remark}
Definition~\ref{def:memory} makes ``memory is the trajectory, not the state''
precise at the level of the field: what is remembered is \emph{how the affect
got here}, encoded as the integral of its motion, and the finer $\widetilde M$
separates histories that the coarse endpoint $M$ conflates.
\end{remark}

% ===========================================================================
\section{Metatheory of the Coupled System}
\label{sec:metatheory}
% ===========================================================================

The full system couples a fast allocation (\cref{eq:replicator}) to a slow field
(\cref{eq:homeostasis}), with the field entering the allocation through the
fitness $f_i(\field)$ and the allocation entering the field through the
interoceptive signals $\coh(\alloc),\varepsilon(\alloc)$. We record its
guarantees.

\subsection{Boundedness}

\begin{proposition}[Boundedness]
\label{prop:bounded}
If $\setpt$ is bounded on $\Unit\times\Rp$ and $V$ is bounded on the (compact)
region containing the population, then along any trajectory the allocation stays
in $\simplex^{N-1}$ (\cref{thm:conservation}) and the field stays in a compact
ball: $\limsup_{t}\norm{\field(t)}\le\sup\norm{\setpt}$.
\end{proposition}
\begin{proof}
The simplex is forward-invariant by \cref{thm:conservation}. For the field, with
$B^\ast=\sup\norm{\setpt}$, outside the ball $\norm{\field}\le B^\ast+\delta$ we
have $\frac{\dd}{\dd t}\tfrac12\norm{\field}^2=\inner{\field}{-\gamma(\field-
\setpt)}\le-\gamma\norm{\field}^2+\gamma\norm{\field}B^\ast<0$, so the ball is
attracting; hence $\limsup\norm{\field}\le B^\ast$ \citep{khalil2002}.
\end{proof}

\subsection{The coherence phase transition}

A central qualitative prediction is that the population's coherence undergoes a
sharp transition as the mutual coupling among candidates increases---the same
mean-field transition that organises synchronization
\citep{kuramoto1984,strogatz2000}. We give a mean-field version.

\begin{assumption}[Mean-field coupling]
\label{ass:meanfield}
Suppose the field setpoint reinforces alignment: candidates are rewarded for
agreeing with the population's mean phase, so that the effective fitness
contains a term $K\,\coh(\alloc)\cos(\theta_i-\psi)$ with coupling $K\ge0$,
where $\theta_i$ is candidate $i$'s phase angle, $\psi$ the mean-phase angle,
and the candidates' intrinsic phases are drawn from a symmetric, unimodal
density $\rho(\omega)$ with spread $\sigma_\omega$.
\end{assumption}

\begin{theorem}[Coherence transition]
\label{thm:transition}
Under \cref{ass:meanfield} the fully incoherent allocation
($\coh=0$) is the only stationary coherence for $K<K_c$ and becomes unstable for
$K>K_c$, where the critical coupling is
\begin{equation}
K_c\;=\;\frac{2}{\pi\,\rho(0)}\quad(\text{equivalently } K_c=c\,\sigma_\omega
\text{ for a fixed shape constant }c).
\label{eq:Kc}
\end{equation}
For $K>K_c$ a coherent branch $\coh(K)>0$ emerges continuously,
$\coh(K)\sim\sqrt{K-K_c}$ near onset. Thus the population exhibits a
\emph{turbulent} regime ($\coh\approx0$) below $K_c$ and a \emph{coherent}
regime ($\coh>0$) above it.
\end{theorem}
\begin{proof}[Proof sketch]
This is the mean-field synchronization transition. Writing the self-consistency
condition $\coh=K\coh\int_{-\pi/2}^{\pi/2}\cos^2\theta\,\rho(K\coh\sin\theta)\,
\dd\theta$ and expanding for small $\coh$ gives $1=K\,\tfrac{\pi}{2}\rho(0)+
O(\coh^2)$, whence the onset $K_c=2/(\pi\rho(0))$ and the square-root branch by
the next order \citep{strogatz2000,kuramoto1984}. Below $K_c$ the only solution
is $\coh=0$; above it the incoherent solution is linearly unstable and the
coherent branch is stable.
\end{proof}

\begin{remark}[Regimes as affective states]
\cref{thm:transition} identifies the two coarse affective regimes with the two
phases of a synchronization transition: dispersed/turbulent processing below the
critical coupling, aligned/coherent processing above it. Because the coupling
$K$ is itself set by the field's setpoint (\cref{ass:meanfield}), the system can
drive itself across the transition by interoceptive feedback---deliberately
de-cohering to explore, or cohering to commit.
\end{remark}

\subsection{Stability of the coupled slow--fast system}

\begin{theorem}[Slow--fast stability]
\label{thm:slowfast}
Suppose the allocation is fast and the field slow: in \cref{eq:replicator}
replace time by $t$ and in \cref{eq:homeostasis} by $\eta t$ with
$0<\eta\ll1$. Assume (i) for each fixed $\field$ the fast subsystem
\cref{eq:replicator} has the globally asymptotically stable equilibrium
$\alloc^\beta(\field)$ (\cref{thm:repconv}), depending smoothly on $\field$
(\cref{thm:gibbs}); and (ii) the reduced slow dynamics
$\dot\field=-\gamma(\field-\setpt(\coh(\alloc^\beta(\field)),
\varepsilon(\alloc^\beta(\field))))$ has a locally exponentially stable
equilibrium $\field^\infty$. Then for $\eta$ small enough the full coupled
system has a locally exponentially stable equilibrium near
$(\alloc^\beta(\field^\infty),\field^\infty)$.
\end{theorem}
\begin{proof}[Proof sketch]
This is Tikhonov's theorem for singularly perturbed systems
\citep{kokotovic1986,khalil2002}. The fast allocation relaxes to the slow
manifold $\alloc=\alloc^\beta(\field)$ (an isolated, exponentially stable root
of the fast equation by (i)); on that manifold the field evolves by the reduced
dynamics, exponentially stable by (ii). The geometric singular-perturbation
theorem then gives, for sufficiently small $\eta$, a persistent attracting
invariant manifold $O(\eta)$-close to the slow manifold and an equilibrium
inheriting exponential stability from the reduced and boundary-layer systems.
\end{proof}

\begin{corollary}[The field as a slowly varying controller]
\label{cor:timescale}
Under \cref{thm:slowfast} the allocation always tracks the current Gibbs
allocation for the prevailing field, and the field migrates slowly toward the
interoceptive setpoint. The system is therefore well-described, on the slow
timescale, as a field-controlled allocator whose allocation is, at each instant,
the optimal (free-energy-minimising) response to the field.
\end{corollary}

% ===========================================================================
\section{Interface to a Host Runtime}
\label{sec:interface}
% ===========================================================================

The theory is designed to plug into a computational runtime as the component
that biases which candidate computations proceed. We isolate the interface
abstractly, so that the theory is independent of any particular host.

\begin{definition}[Sentiment interface]
\label{def:interface}
A \emph{host} supplies, at each instant: the candidate population $C$, the
intrinsic costs $V(x_i)$, the affordances $\feat(x_i)$, and the interoceptive
signals it can measure (or the data from which $\coh,\varepsilon$ are computed).
The \emph{sentiment module} returns: (i) the field $\field$ (and its memory
$M$); and (ii) the allocation $\alloc=\alloc^\beta(\field)\in\simplex^{N-1}$, a
priority distribution over candidates. The host uses $\alloc$ to schedule
effort and consults $M$ for the affective history.
\end{definition}

\begin{principle}[Sufficiency of the pair $(\alloc,M)$]
\label{prin:sufficiency}
For the purpose of deciding which computations proceed and with what priority,
the pair $(\alloc,M)$ is a sufficient statistic of the field: the host need not
see $\field$ itself, only the allocation it induces and the memory it has
accumulated. This decouples the sentiment module from the host's internal
representation and makes the module replaceable.
\end{principle}

\begin{remark}[Relation to a scheduler]
A host scheduler that draws the next unit of work according to $\alloc$ is,
by \cref{thm:repconv,cor:timescale}, sampling from the free-energy-minimising
distribution for the current field; raising $\beta$ sharpens it toward the
fittest candidate (exploitation), lowering it flattens toward exploration
(\cref{cor:temperature}). The field thus furnishes a principled,
self-tuning, conserved priority that a host can consume without further
modelling of affect.
\end{remark}

% ===========================================================================
\section{A Worked Example}
\label{sec:example}
% ===========================================================================

Consider three candidate lines of inquiry $x_1,x_2,x_3$ with intrinsic costs
$V=(0.6,0.4,0.5)$ and two affordances, ``novelty'' $g_1$ and ``consolidation''
$g_2$, with $\feat(x_1)=(0.9,0.1)$, $\feat(x_2)=(0.2,0.8)$,
$\feat(x_3)=(0.5,0.5)$. Take inverse temperature $\beta=4$.

\paragraph{A consolidating field.} With $\field=(0.2,1.0)$ (low novelty price,
high consolidation price) the fitnesses are
$f_i=-V_i+\inner{\field}{\feat(x_i)}$: $f_1=-0.6+(0.18+0.10)=-0.32$,
$f_2=-0.4+(0.04+0.80)=0.44$, $f_3=-0.5+(0.10+0.50)=0.10$. The Gibbs allocation
$a_i\propto e^{4f_i}$ is $\alloc\approx(0.06,0.69,0.25)$: the consolidating
candidate $x_2$ dominates.

\paragraph{Interoceptive regime change.} Suppose coherence is high
($\coh\approx0.9$) but error has stopped improving. The setpoint $\setpt$
responds by raising the novelty price and lowering the consolidation price, say
$\field\to(1.0,0.2)$. Now $f_1=-0.6+(0.90+0.02)=0.32$,
$f_2=-0.4+(0.20+0.16)=-0.04$, $f_3=-0.5+(0.50+0.10)=0.10$, and
$\alloc\approx(0.55,0.16,0.29)$: the field has reallocated effort to the
exploratory candidate $x_1$ \emph{without any candidate being selected by
fiat}---only the prices changed (\cref{prin:modulation}), the allocation tracked
the new Gibbs minimiser (\cref{thm:repconv}), and the budget was conserved
throughout (\cref{thm:conservation}). The memory $M=\field_{\text{after}}-
\field_{\text{before}}=(0.8,-0.8)$ records the shift.

% ===========================================================================
\section{Related Work}
\label{sec:related}
% ===========================================================================

The components of the theory each have a substantial pedigree, and locating them
precisely clarifies the contribution.

\paragraph{Dynamical systems and bifurcation.} The modulation law and its regime
changes are standard equilibrium and bifurcation theory
\citep{strogatz2015,hirsch2013}; the stability arguments use Lyapunov methods and
LaSalle invariance \citep{khalil2002}.

\paragraph{Feedback and homeostasis.} The interoception law is a homeostatic
feedback in the cybernetic tradition \citep{wiener1948,cannon1929,ashby1960,
astrom2008}, with the coherence and error as the measured signals an observer
would track \citep{kalman1960}.

\paragraph{Evolutionary dynamics and allocation.} The economy is replicator
dynamics with entropy regularisation \citep{taylor1978,hofbauer1998,
maynardsmith1982}; its fixed point is the Gibbs/softmax allocation, the
maximum-entropy distribution \citep{jaynes1957,cover2006} and the unique
minimiser of a convex free energy \citep{boyd2004}. The reading of the field as
a price on affordances is the Lagrangian-duality view of allocation
\citep{boyd2004,kelly1998,arrow1954}, and the whole apparatus is an instance of
the ``attention economy'' for a resource-bounded processor \citep{simon1971}.

\paragraph{Synchronization.} The coherence order parameter and its phase
transition are the mean-field theory of coupled oscillators
\citep{kuramoto1984,strogatz2000,winfree1980}.

\paragraph{Comparators.} The free-energy formulation is formally adjacent to the
free-energy principle for adaptive systems \citep{friston2010} and to the
entropy-regularised objectives of reinforcement learning \citep{sutton2018},
but our free energy is over an \emph{allocation} across candidate computations,
not over actions or beliefs, and the controlling field is interoceptive. The
interpretation of the field as affect follows the interoceptive theory of
emotion \citep{craig2002,damasio1994,barrett2017}.

The novelty is the synthesis: a single field that (a) modulates a variance
landscape, (b) flows under a conservation law with a pump, (c) is driven by the
system's own coherence and error, and (d) clears a budget by a maximum-entropy
economy---together with the theorems that make this a well-posed, convergent,
conservative controller with a memory and a clean host interface.

% ===========================================================================
\section{Conclusion}
\label{sec:conclusion}
% ===========================================================================

We have given a self-contained mathematical theory of a sentiment field: a
circulating, interoceptive, resource-allocating quantity that modulates which
candidate computations stabilise and proceed. The four laws---modulation,
circulation, interoception, economy---fit together into a coupled slow--fast
dynamical system whose guarantees we proved: candidates are the minima of a
field-tilted variance landscape and move smoothly with the field except at
affective regime changes; the budget is conserved under transport with a pump;
the field is driven homeostatically by the system's own coherence and error;
the allocation is the unique free-energy-minimising Gibbs distribution and is
reached by an entropy-regularised replicator flow, with the field as the price
of the budget; coherence undergoes a sharp transition between turbulent and
coherent regimes at a critical coupling; the coupled system is stable under
timescale separation; and memory is the path integral of the field. The theory
exposes a minimal interface---an allocation and a memory---through which a host
runtime can consume sentiment as a self-tuning, conserved, accountable priority
without modelling affect itself. It is intended both as a specification of the
sentiment component of such a runtime and as a stand-alone account of why a
field, rather than a scalar, is the right control variable for allocating the
effort of an open-ended computation.

% ===========================================================================
\bibliographystyle{plainnat}
\bibliography{references}
% ===========================================================================
\end{document}
